\subsection{Execution Information as an Optimization Target}

\textbf{Leveraging Test Cases and Code Execution}: Another line of work uses test cases and code execution information to improve code generation. Some examples include Speculyzer \citep{key2022speak}, CodeT \citep{chen2022codet}, CodeGen-Test \citep{zhong2022codegen}, Coder-Reviewer reranking \citep{zhang2023coder}, MBR-EXEC \citep{shi2022natural}
TCoT \citep{tian2023test},
Algo \citep{zhang2023algo}, Pangu-Coder2 \citep{shen2023pangu}, LEVER \cite{ni2023lever}, and Self-Play \citep{haluptzok2022language}.
The idea of these works is to generate many programs and many test cases and select which programs and test cases seem correct based on the execution results.
using execution info. Other works use RL-style execution feedback to improve code generation, including CodeRL \citep{le2022coderl}, Reflexion \citep{shinn2023reflexion}, and PG-TD \citep{zhang2023planning}. \citep{chen2023teaching, olausson2023demystifying, madaan2023self, peng2023check, zhang2023self} investigate self-repair, using error messages as feedback for models to improve.

Most relevant to our work, a handful of works examine and improve the execution ability of code LMs. \cite{austin2021program},
Scratchpad \citep{nye2021show}, and CodeExecutor \citep{liu2023code} train code LMs on execution information. Inspired by these works, we briefly touch on two primitive ways to improve performance on our benchmark, chain-of-thought and fine-tuning. Moving forward, we believe our \benchmark could serve as a useful reference point as more techniques are designed to improve code execution abilities.
